\documentclass[11pt, a4paper]{article}

\usepackage[utf8]{inputenc}
\usepackage[margin=1in]{geometry}
\usepackage{hyperref}
\usepackage{xcolor}
\usepackage{listings}
\usepackage{tcolorbox}
\usepackage{booktabs}

% Code block styling
\definecolor{codegray}{rgb}{0.95,0.95,0.95}
\definecolor{codeblue}{rgb}{0.1,0.2,0.8}
\definecolor{codegreen}{rgb}{0.1,0.6,0.1}

\lstset{
    backgroundcolor=\color{codegray},
    basicstyle=\ttfamily\small,
    breaklines=true,
    keywordstyle=\color{codeblue}\bfseries,
    commentstyle=\color{codegreen}\itshape,
    numbers=none,
    frame=single,
    rulecolor=\color{black!20},
    aboveskip=1em,
    belowskip=1em
}

\title{\textbf{A Quick-Start Guide to Fine-Tuning SmolVLA on UCL Remote GPUs}}
\author{Edwin Redhead}
\date{\today}

\begin{document}

\maketitle

\begin{abstract}
This guide documents the exact end-to-end workflow used for a successful SmolVLA run on UCL TSG GPUs with pre-converted datasets (001, 002, 003). It covers scratch/home separation, resilient nohup training, progress monitoring, and reliable retrieval of full checkpoints and logs.
\end{abstract}

\tableofcontents
\vspace{1em}
\hrule
\vspace{2em}

% =============================================================================
\section{The Core Concept: Persistent Storage vs.\ Scratch Space}
% =============================================================================

When you log into a managed UCL GPU machine, your home directory (\texttt{\char`\~/}) is mounted via a network drive with a strict \textbf{10GB quota}. SmolVLA's dependency stack is large, and your robot dataset is larger. The table below shows why a naive install directly into \texttt{\char`\~/} will always fail:

\vspace{0.5em}
\begin{center}
\begin{tabular}{lrr}
\toprule
\textbf{Artefact} & \textbf{Size (approx.)} & \textbf{Cumulative} \\
\midrule
PyTorch (CUDA build)         & $\sim$2.5 GB & 2.5 GB \\
lerobot + SmolVLA deps       & $\sim$1.5 GB & 4.0 GB \\
SmolVLA base model (HF Hub)  & $\sim$4.0 GB & 8.0 GB \\
Raw episode recordings       & $\sim$5$+$ GB & \textbf{Quota exceeded} \\
\bottomrule
\end{tabular}
\end{center}
\vspace{0.5em}

\textbf{The solution:} every GPU workstation provides a dedicated local disk at \texttt{/scratch0/\$USER/} with no quota limit. However, this scratch space is \textbf{wiped entirely} when your booked session ends (72-hour maximum).

The pipeline uses the following split:

\begin{itemize}
    \item \textbf{Code only} stays in the \textbf{persistent home directory} (\texttt{\char`\~/smolvla\_project/}). Source code is small (megabytes) and must survive booking expiry.
    \item \textbf{Everything heavy} goes to \textbf{scratch} (\texttt{/scratch0/\$USER/}): the Python virtual environment, PyTorch, the HuggingFace model cache, raw episode recordings, and the converted LeRobotDataset.
    \item \textbf{Checkpoints} are written to scratch during training (fast NVMe I/O) then automatically \texttt{rsync}'d back to \texttt{\char`\~/smolvla\_project/checkpoints/} when the run finishes, so they survive the booking wipe.
\end{itemize}

\begin{tcolorbox}[colback=red!5!white,colframe=red!75!black,title=CRITICAL: Dataset Lives in Scratch, Not Home]
A common mistake is rsyncing the \texttt{lerobot\_datasets/} or \texttt{raw\_recordings/} directories into your home directory. Even a few episodes of RGB video from a ZED camera will exceed the 10GB quota. \textbf{Never copy large data to \texttt{\char`\~/}.} All data directories must go to \texttt{/scratch0/\$USER/}.
\end{tcolorbox}

% =============================================================================
\section{The tcsh / bash Shell Incompatibility}
% =============================================================================

UCL TSG machines use \textbf{tcsh} as the default interactive shell. All provided shell scripts are written in \textbf{bash} and use \texttt{export}, which is a bash-only keyword that will fail with \texttt{export: Command not found} if executed in tcsh.

\textbf{Always invoke the scripts with an explicit \texttt{bash} prefix:}

\begin{lstlisting}[language=bash]
bash ~/smolvla_project/SmolVLA-Testing/setup_scratch.sh
\end{lstlisting}

Do \textbf{not} run them as:

\begin{lstlisting}[language=bash]
./setup_scratch.sh        # Wrong — runs in tcsh if that is your default shell
source setup_scratch.sh   # Wrong — sources into tcsh, export will fail
\end{lstlisting}

\texttt{setup\_scratch.sh} writes \emph{two} activation shims to scratch at the end of setup:
\begin{itemize}
    \item \texttt{activate\_smolvla.sh} --- for bash (used by all scripts and nohup)
    \item \texttt{activate\_smolvla.csh} --- for interactive tcsh sessions
\end{itemize}

To activate the environment manually in an interactive session:
\begin{lstlisting}[language=bash]
# If your interactive shell is tcsh (UCL default):
source /scratch0/$USER/activate_smolvla.csh

# If you have already switched to bash:
source /scratch0/$USER/activate_smolvla.sh
\end{lstlisting}

% =============================================================================
\section{Step 1: Sync SmolVLA-Testing to Scratch}
% =============================================================================

In this run, \texttt{lerobot} already existed in persistent home (\texttt{\char`\~/smolvla\_project/lerobot}). We synced only \texttt{SmolVLA-Testing} from local to scratch on the GPU node:

\begin{lstlisting}[language=bash]
GPU_NODE="your-gpu-node.cs.ucl.ac.uk"
rsync -avz --progress \
    -e "ssh -J $USER@knuckles.cs.ucl.ac.uk" \
    --exclude 'lerobot_datasets' \
    --exclude 'raw_recordings' \
    --exclude 'checkpoints' \
    /Volumes/ROS2_SSD/SmolVLA-Testing/ \
    $USER@${GPU_NODE}:/scratch0/$USER/SmolVLA-Testing/
\end{lstlisting}

% =============================================================================
\section{Step 2: Transfer Converted LeRobotDataset Directories to Scratch}
% =============================================================================

For this run, datasets were already processed locally into \texttt{lerobot\_datasets/001}, \texttt{002}, and \texttt{003}. We therefore skip raw-recording transfer and conversion, and copy the converted LeRobotDataset directories directly to scratch.

From your \textbf{local machine}, use the jump host to copy directly into scratch on your allocated GPU node:

\begin{lstlisting}[language=bash]
# Set this once for your current booking
GPU_NODE="your-gpu-node.cs.ucl.ac.uk"

# 1) Create destination parent directory on the GPU node
ssh -J $USER@knuckles.cs.ucl.ac.uk $USER@${GPU_NODE} \
    'mkdir -p /scratch0/$USER/lerobot_datasets'

# 2) Copy converted datasets 001/002/003 from local disk to scratch
for DS in 001 002 003; do
    rsync -avzP \
        -e "ssh -J $USER@knuckles.cs.ucl.ac.uk" \
        "/Volumes/ROS2_SSD/SmolVLA-Testing/lerobot_datasets/${DS}/" \
        "$USER@${GPU_NODE}:/scratch0/$USER/lerobot_datasets/${DS}/"
done
\end{lstlisting}

After this, confirm the data arrived on the GPU node:
\begin{lstlisting}[language=bash]
ssh -J $USER@knuckles.cs.ucl.ac.uk $USER@${GPU_NODE} \
    'ls -lah /scratch0/$USER/lerobot_datasets && du -sh /scratch0/$USER/lerobot_datasets/{001,002,003}'
\end{lstlisting}

\begin{tcolorbox}[colback=yellow!10!white,colframe=orange!80!black,title=Non-blocking Remote Noise]
You may see \texttt{VBoxManage} warnings when connecting to the node. In this workflow, those messages were noisy shell startup output and did not block \texttt{ssh} or \texttt{rsync}. The transfer blocker was missing parent directories under \texttt{/scratch0/\$USER/}.
\end{tcolorbox}

% =============================================================================
\section{Step 3: Connect to the GPU Node}
% =============================================================================

Once your TSG GPU booking begins, connect via SSH:

\begin{lstlisting}[language=bash]
GPU_NODE="your-gpu-node.cs.ucl.ac.uk"
ssh -l $USER -J $USER@knuckles.cs.ucl.ac.uk ${GPU_NODE}
\end{lstlisting}

Alternatively use the \textbf{Apache Guacamole} web interface (URL in your booking confirmation). This opens a Rocky Linux graphical desktop in your browser.

Confirm the GPU is allocated and visible:
\begin{lstlisting}[language=bash]
hostname          # Confirm you are on the correct node
nvidia-smi        # Should list the GPU with low utilisation
\end{lstlisting}

% =============================================================================
\section{Step 4: Build the Scratch Environment}
% =============================================================================

Run \texttt{setup\_scratch.sh} from the GPU node. This script:

\begin{enumerate}
    \item Exports all cache environment variables (\texttt{UV\_CACHE\_DIR}, \texttt{HF\_HOME}, etc.) pointing to scratch \textbf{before any installation begins}.
    \item Creates the scratch directory tree.
    \item Verifies \texttt{uv} is on \texttt{PATH} (installs it to scratch if not).
    \item Runs \texttt{uv sync --extra smolvla --extra dataset} inside the \texttt{lerobot} project, creating a full virtual environment at \texttt{/scratch0/\$USER/smolvla\_venv/}.
    \item Verifies PyTorch can see the CUDA GPU.
    \item Writes both a bash and a tcsh activation shim to scratch.
    \item Runs \texttt{patch\_nvenc.py} to fix the NVENC API version mismatch (see Section~\ref{sec:nvenc}).
\end{enumerate}

\begin{lstlisting}[language=bash]
bash /scratch0/$USER/SmolVLA-Testing/setup_scratch.sh
\end{lstlisting}

The setup takes approximately 5--15 minutes. A successful run ends with:
\begin{lstlisting}
    torch version  : 2.x.x+cu121
    CUDA available : True
    CUDA device    : NVIDIA RTX A5000
[5/5] Patching lerobot video_utils.py for h264_nvenc compatibility...
Patched successfully: .../video_utils.py
=================================================
  Setup complete.
=================================================
\end{lstlisting}

\subsection{The NVENC API Version Mismatch}
\label{sec:nvenc}

lerobot uses PyAV 15.x, which bundles ffmpeg 7.x. ffmpeg 7.x requires NVENC API 13.0. On current UCL TSG images, the GPU driver often exposes NVENC API 12.2, causing encoding to fail immediately with:

\begin{lstlisting}
Cannot load nvcuda.dll / Required NVENC API version 13.0, Found 12.2
\end{lstlisting}

Downgrading PyAV is not possible because Python 3.14 (the version on UCL machines) has no pre-built wheels for \texttt{av < 15}.

\textbf{The fix} (\texttt{patch\_nvenc.py}): patches lerobot's \texttt{encode\_video\_frames} function to call the \textit{system} \texttt{ffmpeg} binary (Lavc59 / ffmpeg 5.x, compiled against NVENC SDK 12.x) via \texttt{subprocess} whenever \texttt{vcodec=h264\_nvenc} is requested. The system ffmpeg works correctly with the installed driver. The patch also forces \texttt{-pix\_fmt yuv420p} because the default \texttt{nv12} pixel format is not recognised by lerobot's frame reader.

The patch is applied automatically by \texttt{setup\_scratch.sh}. It is idempotent — safe to re-run if you need to update it.

% =============================================================================
\section{Step 5: Optional Conversion (Skipped in This Run)}
% =============================================================================

This step is only required when you have not already built \texttt{lerobot\_datasets/<id>} locally. In the April 2026 workflow, conversion was skipped because \texttt{001/002/003} were already available and rsync'd directly to scratch in Step 2.

\textbf{Edit \texttt{convert\_all.sh}} if your session IDs differ from \texttt{001 002 003}:

\begin{lstlisting}[language=bash]
# At the top of convert_all.sh, change this line:
DATASETS=(001 002 003)
\end{lstlisting}

Launch the conversion under \texttt{nohup} using \texttt{restart\_conversion.sh}:

\begin{lstlisting}[language=bash]
bash ~/smolvla_project/SmolVLA-Testing/restart_conversion.sh
\end{lstlisting}

This script kills any existing conversion processes, relaunches \texttt{convert\_all.sh} under \texttt{nohup} (so it survives SSH disconnection), and tails the log. Press \texttt{Ctrl+C} to stop watching --- the conversion continues in the background.

\subsection{Monitoring Conversion Progress}

\begin{lstlisting}[language=bash]
# Reattach to the live log at any time
tail -f /scratch0/$USER/conversion.log

# See only episode-level progress lines
grep "Episode\|Done:\|dataset" /scratch0/$USER/conversion.log

# Count completed episodes in a dataset
ls /scratch0/$USER/lerobot_datasets/001/videos/observation.images.ee_zed_m/ \
    | wc -l
\end{lstlisting}

Typical conversion time is 2--5 minutes per episode when using \texttt{h264\_nvenc}. A 25-episode session takes approximately 1--2 hours total.

\subsection{Check Conversion Reports}

After each dataset finishes, \texttt{data\_converter.py} writes a JSON report:

\begin{lstlisting}[language=bash]
cat /scratch0/$USER/lerobot_datasets/001_report.json
\end{lstlisting}

This contains per-episode success/failure status, episode lengths, and any skipped frames.

% =============================================================================
\section{Step 6: Launch the Training Run}
% =============================================================================

With datasets present in scratch, launch training using \texttt{run\_training.sh}. For this run we jointly trained datasets \texttt{001}, \texttt{002}, and \texttt{003}:

\begin{lstlisting}[language=bash]
bash /scratch0/$USER/SmolVLA-Testing/run_training.sh 001 002 003
\end{lstlisting}

Key training configuration in \texttt{run\_training.sh} for future runs:
\begin{itemize}
    \item \texttt{BATCH\_SIZE=8}
    \item \texttt{STEPS=20000}
    \item \texttt{SAVE\_FREQ=1000} (updated from 500)
    \item \texttt{LOG\_FREQ=50}
    \item \texttt{AMP\_FLAG="--use-amp"}
    \item \texttt{ALLOW\_NEAREST\_FRAME\_FALLBACK=true}
    \item \texttt{ALLOW\_MISSING\_TASK\_FALLBACK=true}
\end{itemize}

The script performs pre-flight checks (venv present, dataset valid), then launches \texttt{main.py} under \texttt{nohup}. Outputs:

\begin{lstlisting}
=================================================
  Training launched.
  PID : 12345
=================================================
  Monitor live output:
        tail -f /scratch0/$USER/smolvla_outputs/001_002_003.log
\end{lstlisting}

\textbf{What \texttt{nohup} does:} closing an SSH terminal sends \texttt{SIGHUP} to foreground processes. \texttt{nohup} immunises the training process against this signal.

\textbf{What the EXIT trap does:} registers a bash \texttt{EXIT} trap that fires on normal completion, Python crash, or OOM kill and \texttt{rsync}'s the entire output directory from scratch to \texttt{\char`\~/smolvla\_project/checkpoints/}.

% =============================================================================
\section{Step 7: Verify the Run is Healthy}
% =============================================================================

\subsection{Check the Training Log}

\begin{lstlisting}[language=bash]
tail -f /scratch0/$USER/smolvla_outputs/001_002_003.log
\end{lstlisting}

A healthy run prints a line every \texttt{--log-freq} steps (default: 50). Look for decreasing loss values. Press \texttt{Ctrl+C} to detach without stopping training.

\subsection{One-Command Status Snapshot}

Use this helper to print current step, percent complete, optional ETA, log file path, and a quick GPU summary:

If your remote \texttt{/scratch0/\$USER/SmolVLA-Testing/} copy predates this helper, sync it once from local first:

\begin{lstlisting}[language=bash]
GPU_NODE="your-gpu-node.cs.ucl.ac.uk"
rsync -avzP \
    -e "ssh -J $USER@knuckles.cs.ucl.ac.uk" \
    /Volumes/ROS2_SSD/SmolVLA-Testing/training_status_snapshot.sh \
    $USER@${GPU_NODE}:/scratch0/$USER/SmolVLA-Testing/
\end{lstlisting}

\begin{lstlisting}[language=bash]
bash /scratch0/$USER/SmolVLA-Testing/training_status_snapshot.sh 001_002_003 20000
\end{lstlisting}

If you keep \texttt{STEPS=20000} in \texttt{run\_training.sh}, the second argument can be omitted.

To estimate completion from the latest saved checkpoint (\texttt{STEPS=20000}, \texttt{SAVE\_FREQ=1000}):
\begin{lstlisting}[language=bash]
latest=$(ls -1 /scratch0/$USER/smolvla_outputs/001_002_003_smolvla/checkpoints/ \
    | grep -E '^[0-9]+$' | sort -n | tail -1)
echo "Latest step: ${latest}/20000"
awk -v s="$latest" 'BEGIN { printf("Progress: %.1f%%\\n", (s/20000)*100) }'
\end{lstlisting}

\begin{tcolorbox}[colback=red!5!white,colframe=red!75!black,title=Signs of a Silent Failure]
If the log stops updating, check the last 50 lines for a Python traceback:
\begin{lstlisting}[language=bash]
tail -n 50 /scratch0/$USER/smolvla_outputs/001_002_003.log
\end{lstlisting}
Common causes: CUDA out-of-memory (reduce \texttt{BATCH\_SIZE} in \texttt{run\_training.sh}), a missing dataset path, or a checkpoint path pointing to a wiped scratch directory.
\end{tcolorbox}

\subsection{Check GPU Utilisation}

\begin{lstlisting}[language=bash]
watch -n 2 nvidia-smi
\end{lstlisting}

A healthy run keeps the GPU at 80--100\% utilisation. If utilisation is near 0\% but the process is alive, the dataloader is the bottleneck --- increase \texttt{NUM\_WORKERS} in \texttt{run\_training.sh}.

\subsection{Check Disk Quota}

\begin{lstlisting}[language=bash]
du -sh ~/                   # Must stay under 10GB
df -h /scratch0/$USER/      # Should have plenty of headroom
\end{lstlisting}

% =============================================================================
\section{Step 8: Safely Disconnecting}
% =============================================================================

\begin{tcolorbox}[colback=yellow!10!white,colframe=orange!80!black,title=CRITICAL: Do Not Click ``Log Out'']
If you are connected via the \textbf{Apache Guacamole} web interface, \textbf{DO NOT} click ``Log Out'' in the system menus. This sends \texttt{SIGTERM} to your entire session, force-killing both the conversion and training processes.

Simply \textbf{close the browser tab}. The graphical session remains active. Your jobs continue unaffected. The same applies to SSH: close the terminal window directly rather than typing \texttt{logout} or \texttt{Ctrl+D}.
\end{tcolorbox}

Reattach at any time by opening a new SSH session:
\begin{lstlisting}[language=bash]
tail -f /scratch0/$USER/smolvla_outputs/001_002_003.log
\end{lstlisting}

% =============================================================================
\section{Step 9: Retrieving Checkpoints}
% =============================================================================

The run finished at step \texttt{020000}, but the automatic rescue to home can be incomplete for very large outputs. The reliable end-to-end retrieval path used here was to pull directly from scratch with resumable \texttt{rsync}.

From a \textbf{new terminal on your local machine}:

\begin{lstlisting}[language=bash]
GPU_NODE="your-gpu-node.cs.ucl.ac.uk"
mkdir -p /Volumes/ROS2_SSD/SmolVLA-Testing/checkpoints/001_002_003_smolvla_full

# macOS system rsync is 2.6.9, so use --partial + --inplace (no --append-verify)
rsync -avP \
    --partial \
    --inplace \
    --timeout=120 \
    -e "ssh -o ServerAliveInterval=30 -o ServerAliveCountMax=6 -J $USER@knuckles.cs.ucl.ac.uk" \
    $USER@${GPU_NODE}:/scratch0/$USER/smolvla_outputs/001_002_003_smolvla/ \
    /Volumes/ROS2_SSD/SmolVLA-Testing/checkpoints/001_002_003_smolvla_full/

# Pull launcher log as well
mkdir -p /Volumes/ROS2_SSD/SmolVLA-Testing/checkpoints/001_002_003_smolvla_full/logs
rsync -avP \
    -e "ssh -J $USER@knuckles.cs.ucl.ac.uk" \
    $USER@${GPU_NODE}:/scratch0/$USER/smolvla_outputs/001_002_003.log \
    /Volumes/ROS2_SSD/SmolVLA-Testing/checkpoints/001_002_003_smolvla_full/logs/
\end{lstlisting}

If transfer drops, run the same command again; it resumes.

Validate final checkpoint coverage locally:
\begin{lstlisting}[language=bash]
ls -1 /Volumes/ROS2_SSD/SmolVLA-Testing/checkpoints/001_002_003_smolvla_full/checkpoints/ \
    | sort -n | tail -n 10
# Expect tail to include 020000
\end{lstlisting}

The directory structure:
\begin{lstlisting}
checkpoints/001_002_003_smolvla_full/
    checkpoints/
            001000/
            002000/
      ...
            020000/    <- final checkpoint
    logs/
            001_002_003.log
\end{lstlisting}

% =============================================================================
\section{End State}
% =============================================================================

This workflow completed with final checkpoint \texttt{020000}. For future runs, keep \texttt{SAVE\_FREQ=1000} in \texttt{run\_training.sh} unless you explicitly need denser checkpointing.

\end{document}
